% SOURCE_STATEMENT: "For $\\operatorname{SU}(N)$-monopoles with maximal symmetry breaking, establish the quasi-isometric asymptotics of the monopole metric on the analogous compactification and determine whether they agree qualitatively with the $\\operatorname{SU}(2)$ case."
% Preparation and final self-audit: CONTRIBUTING.md prompts C and D; source/audit artifacts retained in the local physics_statements_2026_10_11 batch.
% !TeX program = lualatex
\documentclass[11pt]{article}
\usepackage[margin=25mm]{geometry}
\usepackage{fontspec}
\setmainfont{Latin Modern Roman}
\usepackage{amsmath,amssymb,mathtools}
\usepackage{unicode-math}
\setmathfont{Latin Modern Math}
\usepackage{xurl,hyperref}
\hypersetup{colorlinks=true,urlcolor=blue}
\setcounter{secnumdepth}{0}
\setlength{\parindent}{0pt}
\setlength{\parskip}{5pt}
\setlength{\emergencystretch}{3em}
\newcommand{\sref}[2]{\hyperlink{source:#1:#2}{[S#2]}}
\newcommand{\cataloguescope}[1]{\paragraph{Recorded target.}#1}
\newcommand{\R}{\mathbb R}
\newcommand{\C}{\mathbb C}
\newcommand{\N}{\mathbb N}
\newcommand{\Z}{\mathbb Z}
\newcommand{\Q}{\mathbb Q}
\newcommand{\F}{\mathbb F}
\newcommand{\eps}{\varepsilon}
\begin{document}
% UnsolvedMath ID 30002074; source catalogue code OWR-11787-005
\setcounter{section}{30002073}
\section{Asymptotic Metrics for Higher-Rank Monopoles}\label{30002074}
{\small UnsolvedMath ID 30002074 · Mathematical Physics}
\subsection{Definitions and mathematical statement}

\paragraph{Shared notation.}$\mathbb N=\{1,2,\ldots\}$; the standard number fields have their usual meanings. All additional conventions are stated locally.


\paragraph{Definitions and conventions.} A finite-energy monopole on $\mathbb R^3$ consists of a connection $A$ and Higgs field $\Phi$ for the gauge group, in the Bogomolny/monopole setting of the source. Framed moduli divide by gauge transformations tending to the identity at infinity. Magnetic charge is determined by the asymptotic Higgs field: one positive integer $k$ for $SU(2)$, or a vector of $N-1$ charges for maximally broken $SU(N)$. The moduli metric is the $L^2$ inner product on infinitesimal field variations perpendicular to gauge directions. Maximal symmetry breaking means the limiting Higgs value is regular, so its centralizer is a maximal torus. The source compactification adds boundaries describing clusters of lower-charge monopoles separating to infinity. Quasi-isometric metric comparison on an asymptotic region means two metrics bound one another by positive constants there, with the covering and cluster-coordinate conventions retained.

\paragraph{Source scope and qualifications.} The target includes arbitrary permitted magnetic-charge vectors and the compactification's cluster strata, rather than a single separated-charge example. The source calls this work in progress; local metric approximations alone do not yield the requested uniform quasi-isometric control across all relevant regions.

\paragraph{Catalogue mathematical target.}
\paragraph{Catalogue wording retained for traceability.}
For $\operatorname{SU}(N)$-monopoles with maximal symmetry breaking, establish the quasi-isometric asymptotics of the monopole metric on the analogous compactification and determine whether they agree qualitatively with the $\operatorname{SU}(2)$ case.

\paragraph{Saved literature context.} The supplied catalogue's assessment dated 2026-08-21 records: Asymptotic metrics for maximally broken SU(N) monopoles are known, but the proposed compactification and the claimed global quasi-isometric control for arbitrary charges remain unresolved; the record itself describes this as work in progress. The assessment is retained as dated provenance; it is not a new comprehensive literature or open-status certification.
\subsection{Short English statement}
Control the higher-rank monopole metric near every compactification boundary and compare it with the SU(2) cluster asymptotics.
\subsection{Sources}
\begin{itemize}
\item[\hypertarget{source:30002074:1}{S1}] ULAM.AI and the UnsolvedMath Contributors, UnsolvedMath, record 30002074 (OWR-11787-005). Catalogue curation under CC BY 4.0; definitions and English exposition added locally.
\url{https://huggingface.co/datasets/ulamai/UnsolvedMath}.
\item[\hypertarget{source:30002074:2}{S2}] Analysis and Geometric Singularities (2012), Oberwolfach Reports, DOI 10.4171/owr/2012/25, pp. 1487--1562; official report source and mathematical context.
\url{https://doi.org/10.4171/owr/2012/25}.
\item[\hypertarget{source:30002074:3}{S3}] Official EMS Press PDF of the same Oberwolfach report; PDF page 21, printed page 1507.
\url{https://ems.press/content/serial-article-files/46396}.
\item[\hypertarget{source:30002074:4}{S4}] Bielawski, Asymptotic metrics for SU(N)-monopoles with maximal symmetry breaking; reference in the saved dated catalogue literature review.
\url{https://arxiv.org/abs/hep-th/9801092}.
\item[\hypertarget{source:30002074:5}{S5}] Kottke--Singer, Partial compactification of monopoles and metric asymptotics; reference in the saved dated catalogue literature review.
\url{https://arxiv.org/abs/1512.02979}.
\end{itemize}
\label{end:30002074}
\end{document}
